\documentclass[11pt,a4paper]{article}
\usepackage[margin=1in]{geometry}
\usepackage{amsmath,amssymb}
\usepackage{booktabs}
\usepackage{enumitem}
\usepackage{hyperref}
\usepackage{xcolor}
\usepackage{fancyhdr}
\usepackage{listings}

\lstset{
  basicstyle=\ttfamily\small,
  breaklines=true,
  frame=single,
  language=Python,
  keywordstyle=\color{blue},
  commentstyle=\color{gray},
  backgroundcolor=\color{gray!5}
}

\pagestyle{fancy}
\fancyhead[L]{Thesis Working Notes}
\fancyhead[R]{March 19, 2026}

\title{%
  \textbf{Reconstruction from Fine-Tuning:\\Catalog of Approaches and Next Steps}\\[0.5em]
  \large Thesis Working Notes — Returning After 3-Week Break
}
\author{}
\date{March 19, 2026}

\begin{document}
\maketitle

%----------------------------------------------------------------------
\section{Context: Where We Left Off}
%----------------------------------------------------------------------

This document picks up the thesis after a \textbf{3-week pause}. Below is a summary of what was done before the break, why we're pivoting, and a catalog of promising approaches for the next phase.

\subsection{Previous Results}

\begin{itemize}
  \item \textbf{Experiment A (KKT on composed weights):} \emph{Failed.}
    The original Haim et al.\ KKT reconstruction was applied to $\theta_{\mathrm{fine}} = \theta_{\mathrm{base}} + BA$, assuming $N=2$ fine-tuning samples.
    Root cause: the composed model satisfies KKT over \emph{all} ${\sim}502$ samples (500 pre-training + 2 fine-tuning).
    The KKT loss plateaued at ${\sim}460$, essentially $\|\theta_{\mathrm{base}}\|^2$ — the pre-training residual that 2 images cannot explain.

  \item \textbf{Experiment B (NTK / linear gradient approximation):} \emph{Succeeded.}
    Used the weight difference $\Delta\theta = \theta_{\mathrm{fine}} - \theta_{\mathrm{base}}$ as a single-step gradient signal:
    \[
      \Delta\theta \;\approx\; -\eta \sum_i c_i \,\nabla_\theta f(\theta_{\mathrm{base}};\, x_i)
    \]
    Achieved SSIM $\sim$0.80--0.99 on MNIST (full model); $\sim$0.80 with LoRA.
    Limitation: only theoretically exact for $T=1$ step.
\end{itemize}

\subsection{Why New Approaches?}

Experiment B works, but it relies on a \textbf{single-step linear approximation}. Real fine-tuning runs for many steps ($T \gg 1$), where the linear assumption breaks down — $\theta$ changes at every step, shifting the gradient landscape. We need methods that handle multi-step fine-tuning and work on harder data (CIFAR-10, not just MNIST).

The original plan was also to try the ``fine-tune + reconstruct from $\Delta\theta$'' idea on CIFAR-10 (32$\times$32 color images), which would be a stronger thesis result than MNIST alone.

%----------------------------------------------------------------------
\section{Problem Statement}
%----------------------------------------------------------------------

\textbf{Given:}
\begin{itemize}
  \item $\theta_{\mathrm{base}}$ — pre-trained model weights (known to the attacker)
  \item $\theta_{\mathrm{fine}}$ — fine-tuned model weights (known to the attacker)
  \item Possibly: learning rate $\eta$, number of steps $T$, labels $y$
\end{itemize}

\textbf{Goal:} Recover the fine-tuning images $x_1, \dots, x_N$.

%----------------------------------------------------------------------
\section{Catalog of Approaches}
%----------------------------------------------------------------------

%--- A ---
\subsection{Approach A: KKT on Absolute Weights (Haim et al.)}

\paragraph{Idea.} The fine-tuned model satisfies KKT stationarity over all training data:
\[
  \theta_{\mathrm{fine}} \;\propto\; \sum_i \lambda_i \, y_i \,\nabla_\theta \Phi(\theta_{\mathrm{fine}};\, x_i)
\]
Run the standard extraction (optimize $x$ and $\lambda$).

\paragraph{Problem.} The sum includes all $N_{\mathrm{base}} + N_{\mathrm{fine}}$ samples. With $N_{\mathrm{base}}=500$, the base data dominates; extraction cannot isolate the fine-tuning signal.

\paragraph{Status.} \colorbox{red!15}{Tested, failed.} KKT loss $\sim$460, reconstructions are noise.

%--- B ---
\subsection{Approach B: NTK / Linear Gradient Approximation}

\paragraph{Idea.} For a single SGD step ($T=1$):
\[
  \Delta\theta = -\eta \sum_i c_i \,\nabla_\theta f(\theta_{\mathrm{base}};\, x_i),
  \qquad c_i = \frac{\sigma\!\bigl(f(\theta_{\mathrm{base}}; x_i)\bigr) - y_i}{N}
\]
Gradients evaluated at $\theta_{\mathrm{base}}$ (known, fixed). Invert to find $x$.

\paragraph{Limitation.} Only exact for $T=1$. For $T>1$, $\theta$ shifts each step and the linearization accumulates error.

\paragraph{Status.} \colorbox{green!15}{Tested, works.} SSIM 0.80--0.99 on MNIST.

%--- C ---
\subsection{Approach C: KKT on $\Delta\theta$ (Speculative)}

\paragraph{Idea.} Treat the weight difference as a standalone model:
\[
  \Delta\theta \;\propto\; \sum_i \lambda_i \, y_i \,\nabla_\theta \Phi(\Delta\theta;\, x_i)
\]
Load $\Delta\theta$ into a \texttt{NeuralNetwork} and run the standard \texttt{extraction.py}.

\paragraph{Why it might work.} If $\theta_{\mathrm{base}}$ is small relative to the update, $\Delta\theta \approx \theta_{\mathrm{fine}}$.

\paragraph{Why it probably won't.} $\Delta\theta$ is not a ``trained model.'' The KKT condition arises from max-margin convergence of gradient descent — there is no reason a weight \emph{difference} should satisfy it.

\paragraph{Status.} \colorbox{yellow!30}{Not tested.} Cheap experiment — worth a quick try.

%--- D ---
\subsection{Approach D: Gradient Inversion on $\Delta\theta$}

\paragraph{Idea.} For small $\eta$, $\Delta\theta \approx -\eta \times \text{gradient}$. Apply gradient inversion methods (Geiping et al., ``Inverting Gradients'') that reconstruct data from observed gradients. These use cosine similarity loss, total variation regularization, and other tricks from the federated learning attack literature.

\paragraph{Differs from B.} Gradient inversion uses domain-specific priors and losses; Experiment B uses raw MSE on weight differences.

\paragraph{Limitations.} Designed for single-batch gradients. Multi-step $\Delta\theta$ is accumulated and nonlinear.

\paragraph{Status.} \colorbox{yellow!30}{Not tested.}

%--- E ---
\subsection{Approach E: Influence Functions}

\paragraph{Idea.} The relationship between fine-tuning data and weight change involves the Hessian:
\[
  \Delta\theta \;\approx\; -H^{-1}\,\nabla_\theta L(\theta_{\mathrm{base}};\, x_{\mathrm{fine}})
\]
Multiply both sides by $H$:
\[
  H\,\Delta\theta \;\approx\; -\nabla_\theta L(\theta_{\mathrm{base}};\, x_{\mathrm{fine}})
\]
This recovers a \emph{curvature-corrected} gradient, then invert.

\paragraph{Differs from B.} Experiment B treats $\Delta\theta$ as a raw gradient. Influence functions correct for the inverse-Hessian scaling of SGD. The diagonal Fisher approximation to $H$ is cheap for MLPs.

\paragraph{Status.} \colorbox{yellow!30}{Not tested.}

%--- F ---
\subsection{Approach F: Output Divergence Search}

\paragraph{Idea.} The fine-tuning data is what the two models \emph{disagree on most}. Optimize directly in input space:
\[
  x^* = \arg\max_x \;\bigl\|f(\theta_{\mathrm{fine}};\, x) - f(\theta_{\mathrm{base}};\, x)\bigr\|^2
\]
No weight-space math. Just find inputs where the model changed its mind.

\paragraph{Pros.} Model-agnostic. No assumptions about $\eta$, $T$, or training procedure.

\paragraph{Cons.} May converge to adversarial-like inputs rather than natural images. Needs regularization (TV norm, image prior).

\paragraph{Status.} \colorbox{yellow!30}{Not tested.} Trivial to implement — good sanity check / initialization for other methods.

%--- G ---
\subsection{Approach G: Differentiable Unrolling (MAML-style) \texorpdfstring{$\bigstar$}{*}}

\paragraph{Idea.} Simulate the entire fine-tuning process \emph{differentiably}. Find $x$ such that $T$ SGD steps on $x$, starting from $\theta_{\mathrm{base}}$, land at $\theta_{\mathrm{fine}}$:
\begin{align}
  \theta_0 &= \theta_{\mathrm{base}} \notag\\
  \theta_{t+1} &= \theta_t - \eta\,\nabla_\theta L(\theta_t;\, x) \qquad \text{for } t = 0,\dots,T{-}1 \label{eq:inner}\\[4pt]
  \mathcal{L}_{\mathrm{outer}} &= \|\theta_T - \theta_{\mathrm{fine}}\|^2 \notag
\end{align}
Backpropagate $\partial\mathcal{L}_{\mathrm{outer}}/\partial x$ through all $T$ steps (using \texttt{create\_graph=True} in PyTorch).

\paragraph{Why \texttt{create\_graph=True}?}
Normal \texttt{autograd.grad} computes gradients but discards the computation graph. With \texttt{create\_graph=True}, PyTorch builds a \emph{graph of the gradient computation itself}, enabling differentiation \emph{through} the SGD steps. This is the same trick used in MAML (Finn et al., 2017).

\paragraph{The computation chain.}
\[
  x \;\longrightarrow\; f(\theta_0, x) \;\longrightarrow\; \nabla L \;\longrightarrow\; \theta_1 \;\longrightarrow\; f(\theta_1, x) \;\longrightarrow\; \nabla L \;\longrightarrow\; \theta_2 \;\longrightarrow\; \cdots \;\longrightarrow\; \theta_T
\]
$x$ feeds into \emph{every} step — both the forward pass and the gradient update.

\paragraph{Pseudocode.}

\begin{lstlisting}
theta_base = load_state_dict("base.pth")
theta_fine = load_state_dict("fine.pth")
eta = 0.01   # must match actual fine-tuning LR
T = 50       # must match actual fine-tuning steps

x = torch.randn(N, 3, 32, 32, requires_grad=True)
y = torch.tensor([...])  # assumed known

opt = torch.optim.Adam([x], lr=0.001)

for outer_step in range(1000):
    # --- Inner loop: simulate T fine-tuning steps ---
    theta = {k: v.clone() for k, v in theta_base.items()}
    for t in range(T):
        out = functional_forward(theta, x)
        loss_inner = F.binary_cross_entropy_with_logits(out, y)
        grads = torch.autograd.grad(
            loss_inner, theta.values(), create_graph=True
        )
        theta = {k: v - eta * g
                 for (k, v), g in zip(theta.items(), grads)}

    # --- Outer loss: did we arrive at theta_fine? ---
    outer_loss = sum(
        ((theta[k] - theta_fine[k]) ** 2).sum()
        for k in theta
    )
    outer_loss.backward()
    opt.step()
    opt.zero_grad()
    x.data.clamp_(-1, 1)
\end{lstlisting}

\paragraph{Properties.}

\begin{center}
\begin{tabular}{ll}
\toprule
\textbf{Property} & \textbf{Detail} \\
\midrule
Exact for $T$ steps & No linearization; full nonlinear trajectory \\
Reduces to B for $T{=}1$ & $\|\theta_1 - \theta_{\mathrm{fine}}\|^2 = \|\Delta\theta + \eta\nabla L(\theta_0; x)\|^2$ \\
Memory & Stores $T$ computation graphs. MLP with $T{\leq}100$: fine. $T{=}10000$: OOM \\
Must know $\eta, T$ & Attacker needs fine-tuning hyperparameters; can sweep if unknown \\
Local minima & Non-convex in $x$; use multiple random restarts \\
Labels & Must be known, or optimized jointly as continuous variables \\
\bottomrule
\end{tabular}
\end{center}

\paragraph{Gradient checkpointing.} For large $T$, use \texttt{torch.utils.checkpoint} to trade compute for memory (recompute intermediate steps during backward pass).

\paragraph{Why this is the most promising.} It is the natural generalization of Experiment B. For $T{=}1$ it is mathematically equivalent. For $T>1$ it captures the full nonlinear training dynamics that Experiment B approximates away. If Experiment B works at $T{=}1$, this should work at least as well — and extend to the multi-step regime.

%--- H ---
\subsection{Approach H: Loss-Surface Probing}

\paragraph{Idea.} Fine-tuning data should have \emph{low} loss on $\theta_{\mathrm{fine}}$ but \emph{higher} loss on $\theta_{\mathrm{base}}$:
\[
  x^* = \arg\min_x \;\bigl[L(\theta_{\mathrm{fine}};\, x) - L(\theta_{\mathrm{base}};\, x)\bigr]
\]
Find inputs where the fine-tuned model improved most. Essentially a differentiable membership inference attack.

\paragraph{Pros.} Simple, principled (membership inference is well-studied).

\paragraph{Cons.} Finds images ``near'' the fine-tuning data, not necessarily exact reconstructions.

\paragraph{Status.} \colorbox{yellow!30}{Not tested.}

%----------------------------------------------------------------------
\section{Summary and Recommended Order}
%----------------------------------------------------------------------

\begin{center}
\begin{tabular}{clccl}
\toprule
\textbf{\#} & \textbf{Approach} & \textbf{Theory} & \textbf{Cost} & \textbf{Status} \\
\midrule
B & NTK (Experiment B)         & Sound ($T{=}1$)   & Low      & \colorbox{green!15}{Done, works} \\
G & \textbf{Differentiable Unrolling} & \textbf{Exact (any $T$)} & \textbf{High} & Not tested \\
F & Output Divergence           & Heuristic         & Very low & Not tested \\
H & Loss-Surface Probing        & Heuristic         & Low      & Not tested \\
E & Influence Functions         & Sound             & Medium   & Not tested \\
D & Gradient Inversion          & Sound ($T{=}1$)   & Medium   & Not tested \\
C & KKT on $\Delta\theta$      & Weak              & Low      & Not tested \\
A & KKT on full model           & Sound ($N{=}502$) & High     & \colorbox{red!15}{Failed} \\
\bottomrule
\end{tabular}
\end{center}

\paragraph{Recommended next steps.}
\begin{enumerate}
  \item \textbf{Quick sanity checks:} Implement F (output divergence) and H (loss-surface probing) — each is $<$50 lines, gives intuition about what's recoverable.
  \item \textbf{Main experiment:} Implement G (differentiable unrolling). Start with $T{=}1$ to verify it matches Experiment B, then increase $T$.
  \item \textbf{CIFAR-10:} Extend the best-performing method from MNIST to CIFAR-10 (32$\times$32 color, vehicles vs.\ animals).
  \item \textbf{Refinement:} Try E (influence functions) as an upgrade to the gradient recovery step in B or G.
\end{enumerate}

\end{document}
